\documentclass[aspectratio=169,11pt]{beamer}

\usetheme{Madrid}
\usecolortheme{default}
\usefonttheme{professionalfonts}
\setbeamertemplate{navigation symbols}{}
\setbeamertemplate{footline}[frame number]

\usepackage{amsmath, amssymb, booktabs}

\title{Synthesis, Frontier Questions, And Student Research Designs}
\subtitle{Turning Health And Population Questions Into Labor-Economics Designs}
\author{Special Topic 6: Labor Market, Health, and Population -- Week 6}
\date{}

\begin{document}

\begin{frame}
  \titlepage
\end{frame}

\begin{frame}{Central question}
\begin{itemize}
    \item What does a credible labor-economics research design on health, population, and work look like?
    \item The capstone task is translation: health or population mechanism $\rightarrow$ labor object $\rightarrow$ timing $\rightarrow$ counterfactual $\rightarrow$ welfare.
    \item A successful project is not a generic health, epidemiology, or demography paper with work as an afterthought.
\end{itemize}
\end{frame}

\begin{frame}{Course synthesis}
\small
\begin{tabular}{p{0.16\linewidth} p{0.34\linewidth} p{0.36\linewidth}}
\toprule
Week & Health/population object & Labor-economics object \\
\midrule
1 & health, risk, measurement & capacity, labor supply, job choice, welfare \\
2 & disability and chronic conditions & onset, accommodations, programs, exit \\
3 & fertility, mortality, caregiving & lifecycle allocation and persistence \\
4 & mental health and stress & productivity, stigma, retention, welfare \\
5 & disease, environment, demography & incidence, migration, firms, adjustment \\
6 & synthesis and frontier designs & research architecture and entry points \\
\bottomrule
\end{tabular}
\vspace{0.3cm}
\begin{itemize}
    \item The shared discipline is to name the labor margin before estimating a health effect.
\end{itemize}
\end{frame}

\begin{frame}{Health/population to labor research architecture}
\small
\begin{tabular}{p{0.25\linewidth} p{0.34\linewidth} p{0.29\linewidth}}
\toprule
Element & Design question & Examples \\
\midrule
Mechanism & What health/population force moves work? & disability onset, care, disease, aging \\
Labor outcome & Which labor margin is central? & hours, job choice, productivity, wages \\
Unit & Who adjusts? & worker, household, firm, place, cohort \\
Timing & When does exposure matter? & onset, cumulative dose, lifecycle event \\
Counterfactual & What would have happened otherwise? & pre-path, threshold, unexposed group \\
Welfare object & What do wages miss? & health, risk, flexibility, insurance \\
\bottomrule
\end{tabular}
\end{frame}

\begin{frame}{Where the literature is}
\small
\begin{tabular}{p{0.30\linewidth} p{0.30\linewidth} p{0.28\linewidth}}
\toprule
Segment & Mature evidence & Still under-measured \\
\midrule
Health and labor supply & measurement and retirement; severe shocks & work capacity and task demand \\
Disability/chronic conditions & onset, earnings, program rules & employer accommodation and retention \\
Lifecycle allocation & fertility timing and child penalties & elder care and older-worker flexibility \\
Mental health & absenteeism and presenteeism & stigma, treatment, workplace design \\
Disease/env./demog. & exposure shocks and incidence & migration, firm adaptation, welfare \\
\bottomrule
\end{tabular}
\end{frame}

\begin{frame}{Frontier directions}
\begin{itemize}
    \item Biomarkers, administrative linkage, and better work-capacity measures.
    \item Mental-health treatment, disclosure, stigma, supervision, and workplace design.
    \item Aging, retirement flexibility, caregiving, bridge jobs, and accommodations.
    \item Climate, disease exposure, migration, immobility, and unequal incidence.
    \item Within-country heterogeneity and global settings with informality and weak insurance.
    \item Firm response: screening, scheduling, accommodations, retention, and technology.
\end{itemize}
\end{frame}

\begin{frame}{Research design choices}
\begin{enumerate}
    \item Start with a labor object: participation, hours, productivity, job quality, retention, migration, firm policy, or welfare.
    \item Define the mechanism: diagnosis, onset, treatment, exposure, care burden, demographic composition, or policy.
    \item Choose the unit: worker, household, firm, establishment, place, cohort, or sector.
    \item State timing: anticipation, onset, cumulative exposure, recovery, or long-run scarring.
    \item Name the counterfactual and the main selection concern.
    \item Explain the welfare or distributional object.
\end{enumerate}
\end{frame}

\begin{frame}{Common failure modes}
\begin{itemize}
    \item Starting from a health condition without a labor-market question.
    \item Treating self-report, diagnosis, disability status, treatment, and work capacity as the same object.
    \item Defining exposure timing after seeing the outcome.
    \item Ignoring household or firm responses when they are the mechanism.
    \item Treating migration, program take-up, exit, or retirement only as attrition.
    \item Reporting earnings or employment as welfare without discussing risk, flexibility, stigma, pain, or amenities.
    \item Claiming broad relevance from one setting without a portable labor mechanism.
\end{itemize}
\end{frame}

\begin{frame}{Research opportunities}
\small
\begin{tabular}{p{0.29\linewidth} p{0.31\linewidth} p{0.28\linewidth}}
\toprule
Open area & Why it matters for labor & Plausible contribution \\
\midrule
Accommodations & firms shape feasible jobs & linked data or policy shocks \\
Mental-health stigma & information affects search and retention & disclosure or treatment design \\
Caregiving/flexibility & aging changes household constraints & schedules and retirement margins \\
Climate/disease incidence & exposure changes work and mobility & timing and place-based designs \\
Health inequality & adjustment capacity differs sharply & within-country heterogeneity \\
\bottomrule
\end{tabular}
\end{frame}

\begin{frame}{Research Lab design}
\begin{itemize}
    \item Primary anchor: Meyer--Mok-style disability onset, earnings, income, consumption, and dynamic labor incidence.
    \item Challenge anchor: Albanesi--Kim-style disease exposure, heterogeneous incidence, and aggregate labor-market adjustment.
    \item Reproduce: recover a synthetic event-study profile around disability onset.
    \item Diagnose: classify capacity loss, program incentives, accommodation, household insurance, firm retention, selection, and welfare.
    \item Transfer: redesign the same architecture for mental health, caregiving, climate, disease exposure, aging, or firm accommodation.
\end{itemize}
\end{frame}

\begin{frame}{Capstone standard}
\begin{itemize}
    \item What labor margin changes?
    \item What health or population mechanism moves it?
    \item Who adjusts: worker, household, firm, place, or cohort?
    \item When does exposure, onset, treatment, or adaptation matter?
    \item What counterfactual makes the design credible?
    \item What welfare or distributional object makes the paper important?
\end{itemize}
\end{frame}

\begin{frame}{Takeaways}
\begin{itemize}
    \item Health and population belong in labor economics when they change work, wages, productivity, allocation, firm behavior, inequality, or welfare.
    \item The frontier is less about adding one more health coefficient and more about measuring mechanisms, timing, and adjustment.
    \item Strong projects are explicit about labor object, mechanism, unit, timing, counterfactual, selection, and welfare.
    \item The field is open because health constraints, family constraints, firm response, and demographic change are increasingly central to how labor markets work.
\end{itemize}
\end{frame}

\end{document}
